\section{SPD Information Compression and Canonical Completion}
\label{sec:spd-realization}

Let \(A\in\R^{d\times m}\) have full column rank.  The associated information
compression is
\begin{equation}
  \pi_A:\SPD^d\to\SPD^m,
  \qquad
  \pi_A(P)=A^\top P A.
  \label{eq:spd-compression}
\end{equation}
For \(R,S\in\SPD^m\), write
\begin{equation}
  R\#_tS
  =R^{1/2}(R^{-1/2}SR^{-1/2})^tR^{1/2},
  \qquad 0\le t\le1,
  \label{eq:airm-geodesic}
\end{equation}
for AIRM geodesic interpolation.
For a reference \(P_0\in\SPD^d\), write
\[
  R_0=A^\top P_0A,
  \qquad
  U_0=P_0^{1/2}AR_0^{-1/2},
  \qquad U_0^\top U_0=I_m.
\]
For \(R\in\SPD^m\), define
\begin{align}
  \Psi_{P_0,A}(R)
  &=P_0^{1/2}
  \left[I+U_0(R_0^{-1/2}RR_0^{-1/2}-I)U_0^\top\right]
  P_0^{1/2}
  \nonumber\\
  &=P_0+P_0AR_0^{-1}(R-R_0)R_0^{-1}A^\top P_0.
  \label{eq:canonical-completion}
\end{align}
This expression is invariant under the visible coordinate change
\((A,R)\mapsto(AB,B^\top RB)\), \(B\in GL(m)\).

\begin{proposition}[A reference is necessary]
\label{prop:reference-necessary}
Let \(W\subsetneq E\) and \(R\in\SPD(W)\).  There is no rule depending only
on \((W,R)\) that selects a full \(P\in\SPD(E)\), restricts to \(R\) on
\(W\), and is invariant under every linear automorphism of \(E\) that fixes
\(W\) pointwise.  Canonical ambient completion therefore requires additional
structure such as \(P_0\).
\end{proposition}

\begin{theorem}[Complete SPD information-compression theorem]
\label{thm:spd-compression}
The map \(\pi_A\) and the section \(\Psi_{P_0,A}\) satisfy the following.
\begin{enumerate}[label=(\Alph*),leftmargin=*]
  \item \emph{Contraction and submetry.}  For all \(P,Q\in\SPD^d\),
  \begin{equation}
    \dai(A^\top PA,A^\top QA)\le\dai(P,Q),
    \label{eq:spd-contraction}
  \end{equation}
  and for every \(r\ge0\),
  \begin{equation}
    \pi_A(\overline B_{\AI}(P_0,r))
    =\overline B_{\AI}(R_0,r).
    \label{eq:spd-ball-submetry}
  \end{equation}
  \item \emph{Unique canonical completion.}  For every \(R\in\SPD^m\),
  \begin{equation}
    \boxed{
    \Psi_{P_0,A}(R)
    =\argminop_{A^\top PA=R}\dai(P_0,P),
    \qquad
    \dai(P_0,\Psi_{P_0,A}(R))=\dai(R_0,R).}
    \label{eq:spd-shortest-completion}
  \end{equation}
  The minimizer is unique.
  \item \emph{Horizontal geometry.}  The section \(\Psi_{P_0,A}\) is an
  isometric totally geodesic embedding.  The map \(\pi_A\) is a Riemannian
  submersion; at \(P\), the inverse of
  \(D\pi_A(P)|_{\mathcal H_P}\) is
  \begin{equation}
    h\longmapsto D\Psi_{P,A}(A^\top PA)[h].
    \label{eq:spd-horizontal-lift}
  \end{equation}
  Hence \(\pi_A\) is split-Hadamard and visible AIRM geodesics have unique
  horizontal ambient lifts.
  \item \emph{Invisible-fiber minimax geometry.}  If \(m<d\), the fiber
  \(\mathcal F_A(R)=\{P:A^\top PA=R\}\) has infinite AIRM diameter.  With
  \(P_\star=\Psi_{P_0,A}(R)\) and
  \[
    \mathcal F_A(R;B)
    =\mathcal F_A(R)\cap\overline B_{\AI}(P_\star,B),
  \]
  one has
  \begin{equation}
    \inf_{Q\in\SPD^d}\sup_{P\in\mathcal F_A(R;B)}\dai(Q,P)=B,
    \label{eq:spd-minimax}
  \end{equation}
  and \(P_\star\) is the unique minimax center.
  \item \emph{Exact radial decision reduction.}  Let \(\mathcal A\) be any
  family of admissible \((P_0,A)\), let \(m(A)\) denote the number of columns
  of \(A\), allowing different visible dimensions,
  let \(\rho\) be nondecreasing, and let \(\varphi\) be any extended-real
  visible functional.  Then
  \begin{align}
    &\inf_{\substack{(P_0,A)\in\mathcal A\\P\in\SPD^d}}
    \left\{\rho(\dai(P_0,P))+\varphi(P_0,A,A^\top PA)\right\}
    \nonumber\\
    &\quad=
    \inf_{\substack{(P_0,A)\in\mathcal A\\R\in\SPD^{m(A)}}}
    \left\{\rho(\dai(A^\top P_0A,R))+\varphi(P_0,A,R)\right\}.
    \label{eq:spd-decision-reduction}
  \end{align}
  Every reduced minimizer has the canonical full lift
  \(P=\Psi_{P_0,A}(R)\).  If \(\rho\) is strictly increasing on its finite
  domain, every finite-valued full minimizer is canonical.
\end{enumerate}
\end{theorem}

The closed form requires no constrained SPD optimization.  A stable
implementation forms \(R_0=A^\top P_0A\), uses Cholesky or symmetric solves
instead of explicit inverses, and applies the resulting rank-\(m\) update in
\eqref{eq:canonical-completion}.  Dense assembly costs
\(O(d^2m+dm^2+m^3)\), while a low-rank representation avoids materializing
the full update.  Numerical sensitivity is governed by the conditioning of
\(R_0\); recovery of the descriptor subspace additionally becomes ill
conditioned as an eigenvalue of the normalized target approaches the
reference value one, as quantified in \eqref{eq:stratified-conditioning}.

The set \(\mathcal F_A(R;B)\) is a local AIRM ambiguity model centered at
the reference-selected completion \(P_\star\).  Thus
\eqref{eq:spd-minimax} is a radius-conditioned nonidentifiability statement:
it does not define a reference-free uncertainty set or a reference-free
minimax center.

\begin{corollary}[Visible update quotient and invisible drift]
\label{cor:visible-update}
Let \(\alpha=Ac\ne0\) be a visible covector and let
\(u_P=-P\alpha\).  For every \(P\) satisfying \(A^\top PA=R\),
\begin{equation}
  A^\top u_P=-Rc.
  \label{eq:visible-update-quotient}
\end{equation}
Thus the class of the update modulo \(\ker A^\top\) is fixed.  Conversely,
for every \(n\in\ker A^\top\), there exists a feasible \(P\) such that
\begin{equation}
  u_P=u_{\Psi_{P_0,A}(R)}+n.
  \label{eq:arbitrary-invisible-drift}
\end{equation}
The canonical completion is the unique AIRM-nearest representative and has no
invisible component in the \(P_0\)-orthogonal splitting.
\end{corollary}

\begin{remark}[What is selected]
The visible target \(R\) is an input to the theorem.  The result selects the
least AIRM deformation of \(P_0\) compatible with that target.  It does not
assert that a first-order oracle uniquely determines \(R\), nor that the
reference \(P_0\) is reference-free.
\end{remark}

The same completion also has an information-theoretic characterization.
For Gaussians it can equivalently be read from the KL chain rule as retaining
the reference conditional law on invisible coordinates.  We include it as a
consequence and alternate characterization of the AIRM completion, rather
than as an independent divergence-projection novelty claim.

\begin{corollary}[Gaussian KL completion]
\label{cor:gaussian-kl}
Viewing \(P\) as the covariance of a zero-mean Gaussian, the problem
\begin{equation}
  \min_{A^\top PA=R}
  D_{\mathrm{KL}}\!\left(N(0,P)\,\|\,N(0,P_0)\right)
  \label{eq:gaussian-kl-problem}
\end{equation}
has the unique minimizer \(\Psi_{P_0,A}(R)\).
\end{corollary}

For a fixed information channel, exact decision reduction yields an explicit
regularized update.  Let \(R_0=A^\top P_0A\), let \(\Rhat\in\SPD^m\) be a
positive data-visible target, and let \(0<\tau<1\).

\begin{corollary}[Prior--data geodesic completion]
\label{cor:prior-data}
The full-SPD objective
\begin{equation}
  (1-\tau)\dai(P_0,P)^2
  +\tau\dai(A^\top PA,\Rhat)^2
  \label{eq:prior-data-objective}
\end{equation}
has the unique minimizer
\begin{equation}
  \boxed{
  P_\tau=\Psi_{P_0,A}(R_0\#_\tau\Rhat).}
  \label{eq:prior-data-solution}
\end{equation}
If \(D=\dai(R_0,\Rhat)\), then
\[
  \dai(P_0,P_\tau)=\tau D,
  \qquad
  \dai(A^\top P_\tau A,\Rhat)=(1-\tau)D,
\]
and the minimum objective value is \(\tau(1-\tau)D^2\).  Perturbing the
visible target from \(R_1\) to \(R_2\) changes the completed solution by at
most \(\tau\dai(R_1,R_2)\).
\end{corollary}

The proofs of
\cref{prop:reference-necessary,thm:spd-compression,cor:visible-update,cor:gaussian-kl,cor:prior-data}
are in
\cref{app:spd-proofs}.  Their common mechanism is that orthogonal compression
contracts the whitened Frobenius norm, while the section
\eqref{eq:canonical-completion} attains equality.
